\documentclass[11pt]{article}
\usepackage[margin=1in]{geometry}
\usepackage{amsmath,amssymb,amsthm}
\usepackage{booktabs}
\usepackage{fancyhdr}
\usepackage{xcolor}
\usepackage[colorlinks=true,linkcolor=red,urlcolor=cyan]{hyperref}

\pagestyle{fancy}
\fancyhf{}
\fancyhead[L]{\small Three Benchmarks, Three Winners}
\fancyhead[R]{\small Tenneson, \today}
\fancyfoot[C]{\thepage}
\renewcommand{\headrulewidth}{0pt}

\theoremstyle{plain}
\newtheorem{theorem}{Theorem}[section]
\newtheorem{proposition}[theorem]{Proposition}
\newtheorem{corollary}[theorem]{Corollary}
\theoremstyle{definition}
\newtheorem{definition}[theorem]{Definition}
\theoremstyle{remark}
\newtheorem{remark}[theorem]{Remark}
\newtheorem{finding}{Finding}

\newcommand{\Sig}{\Sigma}
\newcommand{\D}{\mathsf{D}}

\title{\textbf{Three Benchmarks, Three Winners}\\[4pt]
\large A Measured Comparison of the Predator Provers\\[2pt]
\normalsize Premise Selection, Search Policy, and the Cost of a Node}
\author{Brian Tenneson\\ \small M.S., Applied Data Science\\
\small \texttt{btenneson2301@baypath.edu}}
\date{\today}

\begin{document}
\maketitle

\begin{abstract}
Seven systems have been built under the name Predator. They rank different
objects, over different corpora, in different units, and for most pairs the
question ``which is better'' is not well posed: a premise selector's
\emph{fraction of the library read} and a search policy's \emph{nodes expanded}
are not commensurable quantities, and reporting them in one column would be the
easiest available way to mislead a reader.

Partitioned into three benchmarks on which comparison \emph{is} well posed, the
systems separate cleanly. Over set.mm, Predator\_4 wins premise selection on
every metric and its performance is flat from four thousand statements to fifty
thousand. Over the condensed-detachment fragment, Predator\_5 expands $2.14
\times$ fewer nodes than breadth-first search and---this is the claim that had
never been checked---finishes $1.76 \times$ faster in wall clock, while solving
every target and returning a shortest proof $99\%$ of the time.

Two results cut across the benchmarks. A linear ranker beat a random forest in
three independent settings, which is evidence about the features rather than
about the forest. And every search strategy that \emph{discards} states---beam
search, top-$k$ pruning---finished slower in seconds than unguided breadth-first
search \emph{and} solved fewer targets. Proof-covering theory predicts that
discarding forfeits a completeness theorem; these measurements say it also
forfeits speed, so on this fragment the guarantee was free.

A general rule falls out of the timing data, and it is a statement about domains
rather than about policies: a learned policy pays only when expanding a node is
expensive relative to scoring one. Predator\_1's beam search lost $2.5 \times$
in wall clock while opening $18 \times$ fewer nodes, not because its policy was
poor but because its nodes cost $21\,\mu s$ and were not worth avoiding.
\end{abstract}

\tableofcontents

\section{What can be compared}

\begin{center}
\small
\begin{tabular}{llll}
\toprule
\textbf{System} & \textbf{Ranks} & \textbf{Corpus} & \textbf{Metric} \\
\midrule
Predator\_1 & proof-graph nodes & synthetic, 789 & expansions, seconds \\
Predator\_2 & lemmas & set.mm, 8k & recall@$k$, effort \\
Predator\_3 & lemmas & set.mm, 2k--32k & recall@$k$, effort \\
Predator\_4 & lemmas & set.mm, 50k & recall@$k$, effort \\
python\_rf & lemmas & set.mm, 16k & recall@$k$, effort \\
Predator\_5 & edges & CD fragment & expansions vs.\ BFS \\
Predator\_6 & edges & CD fragment & expansions vs.\ BFS \\
\bottomrule
\end{tabular}
\end{center}

Predator\_4's $3.5\times$ is a fraction of the library that must be read.
Predator\_5's $1.8\times$ is nodes expanded before a target appears. These are
different quantities and there is no exchange rate between them. The systems are
therefore judged in three leagues: premise selection over set.mm, search policy
over the fragment, and search strategy timed in seconds.

\section{Theory: what a policy may and may not do}

\begin{definition}[Proof-covering]
A nondeterministic proof-search policy $\Sig$ is \emph{proof-covering} if for
every finite proof $\pi = \langle \pi(1),\dots,\pi(n) \rangle$ from $\Gamma$
there is a branch $b = \langle p_1, p_2, \dots \rangle$ of the search with
$\{\pi(1),\dots,\pi(j)\} \subseteq p_j$ for every $j \le n$.
\end{definition}

Under that hypothesis the Branch-Covering Theorem applies: if $\Gamma \vdash_F
\varphi$ then some target-search branch reaches $\varphi$ by stage
$\|\varphi\|_F(\Gamma)$. Two ways of consuming a learned score have different
standing with respect to it.

\begin{proposition}[Reordering preserves proof-covering]
\label{prop:reorder}
Let $\Sig_{\mathrm{full}}$ return every admissible one-step extension at $p$,
and let $\Sig_\theta$ return the same set permuted by a scoring function
$\theta$. If $\Sig_{\mathrm{full}}$ is proof-covering then so is $\Sig_\theta$.
\end{proposition}

\begin{proof}
A branch is determined by which extensions are \emph{present} in each
$\Sig(p_m,m)$, not by the order in which they are listed. Since
$\Sig_\theta(p,m) = \Sig_{\mathrm{full}}(p,m)$ as sets, the two policies generate
the same set of branches, so any branch shadowing $\pi$ under one shadows it
under the other.
\end{proof}

\begin{proposition}[Truncation destroys it]
\label{prop:prune}
Let $\Sig_{\theta,k}$ return only the top $k$ extensions by $\theta$. Then
$\Sig_{\theta,k}$ need not be proof-covering.
\end{proposition}

\begin{proof}
It suffices that some state have its unique proof-continuing extension scored
outside the top $k$. Nothing ties a fitted $\theta$ to geodesic membership, and
\S\ref{sec:leagueC} reports the failure occurring: pruned and beam policies lost
targets that every complete policy solved.
\end{proof}

Beam search is subject to Proposition~\ref{prop:prune} for the same reason:
retaining the best $W$ states per level discards the rest permanently.

\section{Theory: when a policy is worth its cost}

Node counts are hardware-independent, which is exactly why they can mislead---a
node count does not charge for deciding. Let $E$ denote expansions and $c$ cost
per expansion.

\begin{definition}
The \emph{node speedup} is $\sigma = E_{\mathrm{BFS}} / E_{\pi}$ and the
\emph{per-node overhead} is $\omega = c_{\pi} / c_{\mathrm{BFS}}$.
\end{definition}

\begin{proposition}[Break-even]
\label{prop:breakeven}
The wall-clock speedup of policy $\pi$ over breadth-first search is
$\sigma / \omega$. In particular $\pi$ is faster in real time if and only if
$\sigma > \omega$.
\end{proposition}

\begin{proof}
Total time is expansions times cost per expansion, so
\[
\frac{T_{\mathrm{BFS}}}{T_{\pi}}
= \frac{E_{\mathrm{BFS}}\, c_{\mathrm{BFS}}}{E_{\pi}\, c_{\pi}}
= \frac{\sigma}{\omega}. \qedhere
\]
\end{proof}

Trivial, and it is the proposition the project spent the longest not applying.
Predator\_1 reported $\sigma = 18.1$ with $\omega = 41$, hence
$\sigma/\omega = 0.44$: eighteen times fewer nodes and two and a half times
slower.

\section{League A: premise selection over set.mm}

\begin{center}
\small
\begin{tabular}{lrrrrrr}
\toprule
\textbf{System} & \textbf{Corpus} & \textbf{r@10} & \textbf{r@50} &
\textbf{MRR} & \textbf{Effort} & \textbf{vs.\ brute} \\
\midrule
Predator\_2 (early) & 8k & 0.160 & 0.333 & 0.436 & 0.3513 & $2.6\times$ \\
Predator\_2 (later) & 8k & 0.232 & 0.491 & 0.479 & 0.3871 & $2.4\times$ \\
Predator\_2 (poor features) & 8k & 0.059 & 0.147 & 0.159 & 0.7470 & $1.2\times$ \\
Predator\_3 & 32k & 0.275 & 0.600 & 0.668 & 0.6417 & $1.5\times$ \\
\textbf{Predator\_4} & \textbf{50k} & \textbf{0.419} & \textbf{0.730} &
\textbf{0.885} & \textbf{0.2866} & $\mathbf{3.5\times}$ \\
python\_rf (forest) & 16k & 0.163 & 0.497 & 0.567 & 0.7007 & $1.4\times$ \\
\bottomrule
\end{tabular}
\end{center}

\begin{finding}
Predator\_4 wins League A on every metric, at the largest corpus.
\end{finding}

\begin{finding}[Scaling]
Recall@10 as the corpus grows from 2k to 32k statements:
\[
\text{Predator\_3:}\quad 0.326,\; 0.253,\; 0.143,\; 0.228,\; 0.275
\]
\[
\text{Predator\_4:}\quad 0.336,\; 0.412,\; 0.413,\; 0.431,\; 0.414
\]
Predator\_3 swings by a factor of $2.3$ with no mechanism to explain it and is
measuring variance. Predator\_4 is flat to $\pm 0.02$ from 4k onward, and a flat
curve is itself the result: performance holds as the library grows.
\end{finding}

\begin{finding}[Features dominate models]
The same Predator\_2 code scored recall@10 of $0.059$ and $0.232$ depending only
on which features were active---a fourfold swing, larger than any gap between
versions. In the poor run the weight on local co-citation was $-0.94$; in the
good run, $+0.586$. One sign flip cost more than three versions of progress.
\end{finding}

\section{League B: search policy over the fragment}

The fragment is condensed detachment over implicational formulas: one rule,
most-general unification with occurs check, $\Gamma = \{K,S,W\}$. Its defining
virtue is that $\Sig(p,m)$ is finite and enumerable, which is what permits
breadth-first search to compute true geodesics and hence to supply labels that
are correct by construction rather than read from human-written proofs.

\begin{center}
\small
\begin{tabular}{lll}
\toprule
\textbf{System} & \textbf{Depth 4} & \textbf{Depth 5} \\
\midrule
BFS ($\lambda = 0$) & 25.6 exp. & 75.5 exp. \\
\textbf{Predator\_5} & $\mathbf{15.5 \pm 2.6}$ ($1.60\times$) &
$\mathbf{42.3 \pm 8.6}$ ($1.78\times$) \\
Predator\_6 ($p=0.7$) & indistinguishable & indistinguishable \\
Predator\_6.1 ($p=0.9$) & indistinguishable & --- \\
\bottomrule
\end{tabular}
\end{center}

\begin{finding}[Expert iteration is null here]
Over a $p \times n$ grid paired on identical seeds, no setting beats $n=1$.
Predator\_6.1 came in $0.5$ expansions faster against a seed-to-seed standard
deviation of $2.6$, and Predator\_6 costs $30$--$139\%$ more expansions to reach
the same place.
\end{finding}

The mechanism is clear and it is not a defect in the implementation. Enlarging
the bootstrap frontier creates nothing past the breadth-first horizon, and on
this fragment breadth-first search prices \emph{every} target. Certified labels
are already geodesics; bootstrapped labels are upper bounds, and adding them to
a training set that did not need them lowered the certified share to $22\%$ at
depth 5 and $10\%$ at depth 6.

\begin{remark}[A retracted hypothesis]
An earlier draft claimed expert iteration pays wherever pass one leaves frontier
targets unsolved, inferred from a single depth-5 versus depth-6 comparison. A
controlled sweep refutes it. Depth 6 left $35\%$ of its frontier unsolved---the
most favourable case the hypothesis admits---and Predator\_6 lost there, solving
22 against 24. Depth 4 left $0\%$ and Predator\_6 nominally won. The correlation
runs backwards, and all three differences lay inside the seed noise in any case.
\end{remark}

\section{League C: strategies timed in seconds}
\label{sec:leagueC}

Depth 5, 30 held-out targets, mean over 5 seeds. One ranker, one split; only the
search strategy varies.

\begin{center}
\small
\begin{tabular}{lrrrrrl}
\toprule
\textbf{Strategy} & \textbf{Solved} & \textbf{Exp.} & \textbf{Sec.} &
$\boldsymbol{\mu s}$\textbf{/exp} & \textbf{Opt.} & \textbf{Guarantee} \\
\midrule
BFS ($\lambda=0$) & $100\%$ & $90.0$ & $2.92$ & $1083$ & $100\%$ & geodesic, complete \\
\textbf{A* reorder} & $\mathbf{100\%}$ & $\mathbf{42.0}$ & $\mathbf{1.66}$ &
$1328$ & $\mathbf{99\%}$ & complete \\
A* prune $k{=}4$ & $91\%$ & $77.8$ & $3.48$ & $1487$ & $91\%$ & \textbf{incomplete} \\
Beam $W{=}8$ & $70\%$ & $114.5$ & $19.09$ & $5483$ & $80\%$ & \textbf{incomplete} \\
Beam $W{=}1$ & $35\%$ & $29.4$ & $4.02$ & $4564$ & $41\%$ & \textbf{incomplete} \\
\bottomrule
\end{tabular}
\end{center}

Applying Proposition~\ref{prop:breakeven}:

\begin{center}
\small
\begin{tabular}{lrrr}
\toprule
\textbf{Strategy} & $\boldsymbol{\sigma}$ & $\boldsymbol{\omega}$ &
$\boldsymbol{\sigma/\omega}$ \\
\midrule
\textbf{A* reorder} & $2.14$ & $1.2$ & $\mathbf{1.76}$ \\
A* prune $k{=}4$ & $1.16$ & $1.4$ & $0.84$ \\
Beam $W{=}8$ & $0.79$ & $5.1$ & $0.15$ \\
Beam $W{=}1$ & $3.06$ & $4.2$ & $0.73$ \\
\bottomrule
\end{tabular}
\end{center}

\begin{finding}[The node advantage survives timing]
Predator\_5 achieves $\sigma = 2.14$ at $\omega = 1.2$, hence
$\sigma/\omega = 1.76$: it is faster in seconds, not merely in nodes, while
solving every target and returning a geodesic $99\%$ of the time.
\end{finding}

\begin{finding}[Beam search is the weakest strategy measured]
Beam search---Predator\_1's strategy---solved $70\%$ of targets, expanded
\emph{more} nodes than unguided breadth-first search ($114.5$ against $90.0$),
and finished $6.6\times$ slower in seconds. Its $\omega = 5.1$ is structural: a
level-synchronous beam scores every child of every state in a level, sorts them,
and discards most, paying for scoring it never uses.
\end{finding}

\begin{finding}[Discarding costs more than a theorem]
All three strategies that discard states finished slower than breadth-first
search \emph{and} solved fewer targets. The only strategy that merely reorders
won on every axis. Proposition~\ref{prop:prune} says discarding forfeits the
Branch-Covering Theorem; these measurements say it forfeits time and solve rate
as well. On this fragment there was no speed-for-guarantee trade to make.
\end{finding}

\begin{finding}[Linear beats forest, three times]
python\_rf on set.mm: effort $0.2898 \to 0.7007$ at $6.4\times$ the fit time.
Predator\_3's forest variant: recall@10 $0.185, 0.063, 0.077$. Predator\_5's
ranker: $47.0$ against $34.9$ expansions. Three independent settings agree,
which makes this evidence that the signal is close to linear \emph{in these
features}---not evidence that forests are unsuitable.
\end{finding}

\section{Why the same architecture won here and lost before}

Predator\_1 opened $18\times$ fewer nodes and lost $2.5\times$ in wall clock;
Predator\_5 opens $2.14\times$ fewer and wins $1.76\times$. The smaller node
advantage produced the better outcome. The explanation is in $\omega$, and it is
a statement about the domain.

Breadth-first search costs $1083\,\mu s$ per expansion on the fragment;
Predator\_1's brute force costs $21$--$28\,\mu s$. The difference is not the
logic---both systems close forward from three axioms---but \emph{when the graph
is built}.

Predator\_1 constructs its corpus once, by brute-force closure, reporting stage
sizes $|D_F^{(m-1)}(G)| = [3, 39, 111, 381, 579, 789]$. Search then runs over a
\emph{precomputed} graph, so expanding a node is an adjacency lookup. Scoring it
costs more than visiting it, and $\omega = 41$.

Predator\_5 computes $\Sig(p,m)$ \emph{on demand}, attempting $\D(x,y)$ for every
ordered pair in the state with unification and an occurs check at every
expansion. Visiting a node costs far more than scoring it, and $\omega = 1.2$.

\begin{corollary}
A learned policy is worth its cost precisely when expanding a node is expensive
relative to scoring one---equivalently, when the search graph is generated on
demand rather than precomputed.
\end{corollary}

Predator\_1's policy may have been perfectly good; its nodes were too cheap to be
worth avoiding. The consequence for set.mm is immediate, and it now depends on an
architectural choice rather than on model quality. A prover that searches the
recorded citation DAG has precomputed edges, hence Predator\_1's regime, hence
large $\omega$. A prover that re-derives admissible extensions by unification has
Predator\_5's regime and small $\omega$---but that is the enumeration problem
that drove Predator\_2 to premise selection in the first place.

The two regimes are therefore not freely choosable. Cheap nodes come with a
precomputed graph, which means the proofs are already known; expensive nodes come
with genuine search. This is the central tension of the project stated in units
of microseconds.

\section{Limitations}

The fragment is condensed detachment, not ZFC. Predator\_5's features are stated
in its terms---unifier growth, membership in $\Gamma$---and there is no
$\in$ or $\forall$ for them to describe, so League B and C results do not
transfer to set.mm as stated.

Held-out sets are small: 18 to 39 targets, with seed-to-seed standard deviations
of $2.6$ to $11.7$ expansions. Every figure quoted here is a mean over five or
ten seeds for that reason, and single-run figures from earlier drafts of this
work were not reliable.

League A and League B measure different things and their winners do not compete.
Predator\_4 cannot prove anything: it emits an ordering and stops. Predator\_5
produces proofs but only within a 287-state fragment.

\section{Conclusion}

Three winners, one per benchmark. Predator\_4 for premise selection over set.mm,
on every metric, with performance flat to fifty thousand statements. Predator\_5
for search policy, now the best-evidenced system in the project: fewer nodes,
fewer seconds, complete, and geodesic $99\%$ of the time. Logistic regression
over random forests, in three independent settings.

One loser worth naming: every strategy that discards states. The theory
predicted the completeness cost; the measurements added that there was no
compensating speed. Where reordering suffices, the guarantee is free.

And one rule that outlives the numbers. Whether a learned policy pays is decided
by the ratio of node cost to scoring cost, not by the quality of the policy.
That ratio is a property of the domain, it is measurable before any model is
fitted, and it should be measured first.

\end{document}
